% --- Кодировки и языки ---
\usepackage[T2A]{fontenc}
\usepackage[utf8]{inputenc}
\usepackage[english,russian]{babel}
% --- Страница и колонки ---
% Поля: верх/низ 2.5 см, лево/право 1.75 см.
% Колонки: 8.5 см, промежуток 0.5 см (под A4).
\usepackage[
a4paper,
top=2.5cm,bottom=2.5cm,
left=1.75cm,right=1.75cm,
columnsep=0.5cm
]{geometry}
% --- Шрифт (Times-подобный с кириллицей) ---
% Если редакция требует строго системный Times New Roman — тогда лучше XeLaTeX + fontspec.
\usepackage{tempora}
\usepackage{newtxmath}
% --- Абзацы/переносы/интервалы ---
\usepackage{setspace}
\singlespacing
\usepackage{microtype}
\setlength{\parindent}{1cm} % красная строка 1 см
\setlength{\parskip}{0pt}   % без доп. интервала между абзацами
\usepackage{indentfirst}     % красная строка в первом абзаце после заголовка
% Чтобы LaTeX не «раздувал» вертикальные расстояния ради выравнивания низа колонок:
\raggedbottom
% --- Формулы: размеры 11/7/5 ---
\DeclareMathSizes{11}{11}{7}{5}
% --- Подписи рисунков и таблиц: 10 pt, bold, одинарный ---
\usepackage[font={small,bf},labelsep=period]{caption}
\addto\captionsrussian{\renewcommand\figurename{Рис.}}
\addto\captionsrussian{\renewcommand\tablename{Таблица}}
% --- Графика ---
\usepackage{graphicx}
% --- Ссылки в квадратных скобках [1], порядок по цитированию ---
\usepackage{cite}
% --- Удобные макросы для шапки ---
\newcommand{\udc}[1]{\noindent\textbf{УДК #1}\par\vspace{0.75em}}
\newcommand{\ruTitle}[1]{\begin{center}\textbf{\Large #1}\end{center}\vspace{0.25em}}
\newcommand{\enTitle}[1]{\begin{center}\textbf{\Large #1}\end{center}\vspace{0.25em}}
\newcommand{\ruAuthors}[1]{\begin{center}\textbf{#1}\end{center}\vspace{0.5em}}
\newcommand{\enAuthors}[1]{\begin{center}\textbf{#1}\end{center}\vspace{0.5em}}
\newcommand{\ruAbstract}[1]{\noindent\textbf{Резюме.} \textit{#1}\par\vspace{0.5em}}
\newcommand{\enAbstract}[1]{\noindent\textbf{Abstract.} \textit{#1}\par\vspace{0.5em}}
\newcommand{\ruKeywords}[1]{\noindent\textbf{Ключевые слова:} #1\par\vspace{0.75em}}
\newcommand{\enKeywords}[1]{\noindent\textbf{Keywords:} #1\par\vspace{1.25em}}

\usepackage{titlesec}

\usepackage{titlesec}

\usepackage{titlesec}

\titleformat{\section}
{\bfseries\normalsize}
{}
{0pt}
{}

\titlespacing*{\section}
{0pt}
{\baselineskip}
{0pt}

\titleformat{\subsection}
{\bfseries\normalsize}
{}
{0pt}
{}

\titlespacing*{\subsection}
{0pt}
{0.75\baselineskip}
{0pt}